% --- IMPORTS ---
\documentclass[leqno]{article}
\usepackage{graphicx}
\usepackage{dsfont}
\usepackage{mathtools}
\usepackage{amssymb}
\usepackage{amsmath}
\usepackage{amsthm}
\usepackage{float}
\usepackage{ragged2e}
\usepackage{tensor}
\usepackage{fancyhdr}
\usepackage{empheq}
\usepackage[most]{tcolorbox}
\usepackage{enumitem}
\usepackage{dirtytalk}
\usepackage{marginnote}
\usepackage{microtype}
\usepackage[
  a4paper,
  left=0.8in,
  right=1in,
  top=1in,
  bottom=0.5in,
  headheight=14pt,
  headsep=0.4in
]{geometry}
\setlength{\parindent}{0cm}
% --- TikZ Packages ---
\usepackage{pgfplots}
\pgfplotsset{compat=1.18}
\usepgfplotslibrary{fillbetween, polar}
\usetikzlibrary{calc, positioning}
\usetikzlibrary{angles, quotes}
\usetikzlibrary{arrows.meta}
\usetikzlibrary{decorations.pathreplacing, decorations.markings}
\usetikzlibrary{patterns}
\usetikzlibrary{intersections}
% --- URL Packages ---
\usepackage[colorlinks=true,linkcolor=red,citecolor=magenta,urlcolor=black]{hyperref}%
\urlstyle{same}
% --- END OF IMPORTS ---

% --- CUSTOM COMMANDS ---
\RenewDocumentCommand{\marks}{o m}{%
  \IfValueTF{#1}%
    {\marginnote{\raggedright\textbf{[#2]}}[#1]}%
    {\marginnote{\raggedright\textbf{[#2]}}}%
}
\newcommand{\vspacerule}[2]{%
  \par\vspace*{#1}%
  \noindent\hrulefill\par
  \vspace*{#2}%
}
\definecolor{scarlet}{RGB}{205, 40, 75}
\newcommand{\questionitem}{%
  \item\phantomsection
  \addcontentsline{toc}{section}{Question \number\value{enumi}}%
}
\renewcommand{\contentsname}{Questions}
% --- END OF CUSTOM COMMANDS ---

% --- HEADER CONFIG ---
\pagestyle{fancy}
\fancyhead{}
\fancyfoot{}
\renewcommand{\headrulewidth}{0.9pt}
\lhead{\textit{Solutions} - 1st Order Differential Equations - Integrating Factor}
\chead{}
\rhead{\href{https://danieljsmith.org}{danieljsmith.org}}
% --- END OF HEADER CONFIG ---

\begin{document}
\vspace*{10pt}
\tableofcontents
\newpage
\begin{enumerate}[
  label=\textbf{\arabic*.},
  leftmargin=*,
  topsep=0pt,
  partopsep=0pt,
  parsep=0pt
]

% ---- Question 1 ----
\questionitem

\begin{enumerate}[label=\textbf{(\alph*)}]
  \item Show that an appropriate integrating factor for
  \[
  x^2\,\frac{\mathrm{d}y}{\mathrm{d}x}+y=1,\qquad x>0
  \]
  is $e^{-1/x}$ \marks{4} \\
  \item Hence find the solution of the differential equation
  \[
  x^2\,\frac{\mathrm{d}y}{\mathrm{d}x}+y=1,\qquad x>0
  \]
  for which $y=3$ when $x=1$ \marks{6} \\
\end{enumerate}

\vspace*{10pt}

\begin{tcolorbox}[colback=white, colframe=scarlet, title={\textbf{Solution}}, breakable, grow to left by=6mm, grow to right by=10mm]
\begin{enumerate}[label=\textbf{(\alph*)}]
  \item First write the differential equation in the standard linear form by dividing through by $x^2$:
\[
\frac{\mathrm{d}y}{\mathrm{d}x}+\frac{1}{x^2}y=\frac{1}{x^2}
\]
This is of the form
\[
\frac{\mathrm{d}y}{\mathrm{d}x}+P(x)y=Q(x)
\]
with
\[
P(x)=\frac{1}{x^2}=x^{-2}
\]
For a linear differential equation, an integrating factor is
\[
\mu(x)=e^{\int P(x)\,\mathrm{d}x}
\]
So
\begin{align*}
\mu(x)&=e^{\int x^{-2}\,\mathrm{d}x} \\
&=e^{-x^{-1}} \\
&=e^{-1/x}
\end{align*}
Hence an appropriate integrating factor is $e^{-1/x}$.

  \item Using the integrating factor from part (a), multiply
\[
\frac{\mathrm{d}y}{\mathrm{d}x}+\frac{1}{x^2}y=\frac{1}{x^2}
\]
through by $e^{-1/x}$:
\[
e^{-1/x}\frac{\mathrm{d}y}{\mathrm{d}x}+\frac{e^{-1/x}}{x^2}y=\frac{e^{-1/x}}{x^2}
\]
Now apply the product rule. Since
\[
\frac{\mathrm{d}}{\mathrm{d}x}\left(e^{-1/x}\right)=e^{-1/x}\cdot \frac{1}{x^2}
\]
we get
\begin{align*}
\frac{\mathrm{d}}{\mathrm{d}x}\left(ye^{-1/x}\right)
&=e^{-1/x}\frac{\mathrm{d}y}{\mathrm{d}x}+y\frac{\mathrm{d}}{\mathrm{d}x}\left(e^{-1/x}\right) \\
&=e^{-1/x}\frac{\mathrm{d}y}{\mathrm{d}x}+\frac{e^{-1/x}}{x^2}y
\end{align*}
So the differential equation becomes
\[
\frac{\mathrm{d}}{\mathrm{d}x}\left(ye^{-1/x}\right)=\frac{e^{-1/x}}{x^2}
\]
Integrating both sides with respect to $x$ gives
\[
ye^{-1/x}=\int \frac{e^{-1/x}}{x^2}\,\mathrm{d}x+C
\]
But
\[
\frac{\mathrm{d}}{\mathrm{d}x}\left(e^{-1/x}\right)=\frac{e^{-1/x}}{x^2}
\]
so
\[
\int \frac{e^{-1/x}}{x^2}\,\mathrm{d}x=e^{-1/x}
\]
Hence
\[
ye^{-1/x}=e^{-1/x}+C
\]
Multiplying through by $e^{1/x}$:
\begin{align*}
y&=1+Ce^{1/x}
\end{align*}
Now use the condition $y=3$ when $x=1$:
\begin{align*}
3&=1+Ce \\
2&=Ce \\
C&=\frac{2}{e}
\end{align*}
Substitute back:
\begin{align*}
y&=1+\frac{2}{e}e^{1/x} \\
&=1+2e^{1/x-1}
\end{align*}
Therefore the required solution is $y=1+2e^{1/x-1}$.
\end{enumerate}
\end{tcolorbox}


\newpage

% ---- Question 2 ----
\questionitem

\begin{enumerate}[label=\textbf{(\alph*)}]
  \item Show that
  \[
  \frac{\mathrm{d}}{\mathrm{d}x}\bigl(\ln(\cosh x+1)\bigr)=\tanh\!\left(\frac{x}{2}\right)
  \]
  \marks[-20pt]{3}
  \item Hence solve the differential equation
  \[
  \frac{\mathrm{d}y}{\mathrm{d}x}+\tanh\!\left(\frac{x}{2}\right)y=\sinh x
  \]
  for which $y=2$ when $x=0$. Give your answer in exact form \marks{7} \\
\end{enumerate}

\vspace*{10pt}

\begin{tcolorbox}[colback=white, colframe=scarlet, title={\textbf{Solution}}, breakable, grow to left by=6mm, grow to right by=10mm]
\begin{enumerate}[label=\textbf{(\alph*)}]
  \item Using \(\dfrac{\mathrm{d}}{\mathrm{d}x}\ln u=\dfrac{1}{u}\dfrac{\mathrm{d}u}{\mathrm{d}x}\),
\[
\frac{\mathrm{d}}{\mathrm{d}x}\bigl(\ln(\cosh x+1)\bigr)
=\frac{1}{\cosh x+1}\cdot \frac{\mathrm{d}}{\mathrm{d}x}(\cosh x+1)
=\frac{\sinh x}{\cosh x+1}
\]
Now use the half-angle identities
\[
\sinh x=2\sinh\left(\frac{x}{2}\right)\cosh\left(\frac{x}{2}\right),
\qquad
\cosh x+1=2\cosh^2\left(\frac{x}{2}\right)
\]
So
\begin{align*}
\frac{\sinh x}{\cosh x+1}
&=\frac{2\sinh\left(\frac{x}{2}\right)\cosh\left(\frac{x}{2}\right)}{2\cosh^2\left(\frac{x}{2}\right)} \\
&=\frac{\sinh\left(\frac{x}{2}\right)}{\cosh\left(\frac{x}{2}\right)} \\
&=\tanh\left(\frac{x}{2}\right)
\end{align*}
Hence
\[
\frac{\mathrm{d}}{\mathrm{d}x}\bigl(\ln(\cosh x+1)\bigr)=\tanh\left(\frac{x}{2}\right)
\]
as required.

  \item The differential equation is linear:
\[
\frac{\mathrm{d}y}{\mathrm{d}x}+\tanh\left(\frac{x}{2}\right)y=\sinh x
\]
From part (a),
\[
\frac{\mathrm{d}}{\mathrm{d}x}\bigl(\ln(\cosh x+1)\bigr)=\tanh\left(\frac{x}{2}\right)
\]
so an integrating factor is
\[
\mu(x)=\mathrm{e}^{\int \tanh\left(\frac{x}{2}\right)\,\mathrm{d}x}
=\mathrm{e}^{\ln(\cosh x+1)}
=\cosh x+1
\]
Multiplying the differential equation throughout by \(\cosh x+1\),
\[
(\cosh x+1)\frac{\mathrm{d}y}{\mathrm{d}x}+(\cosh x+1)\tanh\left(\frac{x}{2}\right)y=(\cosh x+1)\sinh x
\]
The left-hand side is
\[
\frac{\mathrm{d}}{\mathrm{d}x}\bigl((\cosh x+1)y\bigr)
\]
Hence the differential equation can be written
\[
\frac{\mathrm{d}}{\mathrm{d}x}\bigl((\cosh x+1)y\bigr)=(\cosh x+1)\sinh x
\]

We can integrate the right hand side by reversing the chain rule:

\[\int (\cosh x+1)\sinh x \,\mathrm{d}x
=\frac{1}{2}(\cosh x+1)^2+C\]

So the differential equation becomes
\begin{align*}
(\cosh x+1)y
&=\int (\cosh x+1)\sinh x \,\mathrm{d}x\\[8pt]
&=\frac{1}{2}(\cosh x+1)^2+C
\end{align*}

Therefore
\[
y=\frac{1}{2}(\cosh x+1)+\frac{C}{\cosh x+1}
\]
Use the condition \(y=2\) when \(x=0\). Since \(\cosh 0=1\),
\begin{align*}
2&=\frac{1}{2}(1+1)+\frac{C}{1+1} \\
2&=1+\frac{C}{2}
\end{align*}
Hence
\[
C=2
\]
Substituting back,
\[
y=\frac{\cosh x+1}{2}+\frac{2}{\cosh x+1}
\]
Therefore the exact solution is
\[
\boxed{
y=\frac{\cosh x+1}{2}+\frac{2}{\cosh x+1}}
\]

Alternatively, one can integrate using the double angle formula for $\sinh 2x$: \[\int (\cosh x + 1)\sinh x \,\mathrm{d}x = \int \cosh x\sinh x + \cosh x \,\mathrm{d}x = \int \frac{1}{2}\sinh 2x + \cosh x \,\mathrm{d}x \]
to obtain an equivalent form of the final answer:

\[
\boxed{
y = \frac{\frac{1}{4}\cosh 2x + \cosh x + \frac{11}{4}}{\cosh x + 1}
}
\]
  \end{enumerate}
\end{tcolorbox}


\newpage

% ---- Question 3 ----
\questionitem

\begin{enumerate}[label=\textbf{(\alph*)}]
  \item For $x>-1$, determine the general solution of the differential equation
  \[
  (1+x)\,\frac{\mathrm{d}y}{\mathrm{d}x}+y=\ln(1+x)
  \]
  giving your answer in the form $y=f(x)$ \marks{4} \\
  \item Given that $y=-1$ when $x=0$, determine the positive value of $x$ for which $y=0$ \marks{2} \\
\end{enumerate}

\vspace*{10pt}

\begin{tcolorbox}[colback=white, colframe=scarlet, title={\textbf{Solution}}, breakable, grow to left by=6mm, grow to right by=10mm]
\begin{enumerate}[label=\textbf{(\alph*)}]
  \item For $x>-1$, divide the differential equation by $1+x$:
\[
\frac{\mathrm{d}y}{\mathrm{d}x}+\frac{1}{1+x}y=\frac{\ln(1+x)}{1+x}
\]

This is a linear differential equation of the form
\[
\frac{\mathrm{d}y}{\mathrm{d}x}+P(x)y=Q(x)
\]
with
\[
P(x)=\frac{1}{1+x}
\]

The integrating factor is
\begin{align*}
\mu(x)&=e^{\int \frac{1}{1+x}\, \mathrm{d}x} \\
&=e^{\ln(1+x)}
\end{align*}

Since $x>-1$, we have $1+x>0$, so
\[
\mu(x)=1+x
\]

Multiplying the differential equation by $1+x$ gives
\begin{align*}
(1+x)\frac{\mathrm{d}y}{\mathrm{d}x}+y&=\ln(1+x)
\end{align*}

The left-hand side is the derivative of $(1+x)y$:
\[
\frac{\mathrm{d}}{\mathrm{d}x}\bigl((1+x)y\bigr)=\ln(1+x)
\]

Integrating both sides with respect to $x$:
\begin{align*}
(1+x)y&=\int \ln(1+x)\, \mathrm{d}x + C
\end{align*}

Now evaluate the integral by letting $u=1+x$, so $\mathrm{d}u=\mathrm{d}x$:
\begin{align*}
\int \ln(1+x)\, \mathrm{d}x&=\int \ln u\, \mathrm{d}u \\
&=u\ln u-u+C \\
&=(1+x)\ln(1+x)-(1+x)+C
\end{align*}

So
\begin{align*}
(1+x)y&=(1+x)\ln(1+x)-(1+x)+C \\
&=(1+x)\ln(1+x)-x-1+C
\end{align*}

Absorbing the constant $-1$ into $C$, we write
\[
(1+x)y=(1+x)\ln(1+x)-x+C
\]

Hence
\begin{align*}
y&=\frac{(1+x)\ln(1+x)-x+C}{1+x} \\
&=\ln(1+x)-\frac{x}{1+x}+\frac{C}{1+x} \\
&=\ln(1+x)-1+\frac{1}{1+x}+\frac{C}{1+x}
\end{align*}

Absorbing the extra constant into $C$ again gives
\[
y=\ln(1+x)-1+\frac{C}{1+x}
\]

Therefore, the general solution is
\[
y=\ln(1+x)-1+\frac{C}{1+x}
\]

  \item Use the condition $y=-1$ when $x=0$.

Substituting into the general solution:
\begin{align*}
-1&=\ln(1+0)-1+\frac{C}{1+0} \\
-1&=0-1+C \\
-1&=-1+C
\end{align*}

So
\[
C=0
\]

Hence
\[
y=\ln(1+x)-1
\]

Now set $y=0$:
\begin{align*}
0&=\ln(1+x)-1 \\
\ln(1+x)&=1
\end{align*}

Exponentiating both sides:
\begin{align*}
1+x&=e \\
x&=e-1
\end{align*}

Since the question asks for the positive value,
\[
x=e-1\approx 1.72
\]

Therefore, the required value is $x=e-1\approx 1.72$.
\end{enumerate}
\end{tcolorbox}


\newpage

% ---- Question 4 ----
\questionitem

\begin{enumerate}[label=\textbf{(\alph*)}]
  \item Show that, for $0<x<\frac{\pi}{2}$, an appropriate integrating factor for
  \[
  \cos x\,\frac{\mathrm{d}y}{\mathrm{d}x}+y\sin x=x
  \]
  is $\sec x$ \marks{4} \\
  \item Hence find the solution of
  \[
  \cos x\,\frac{\mathrm{d}y}{\mathrm{d}x}+y\sin x=x
  \]
  for $0<x<\frac{\pi}{2}$, given that $y=1$ when $x=\frac{\pi}{4}$ \marks{6} \\
\end{enumerate}

\vspace*{10pt}

\begin{tcolorbox}[colback=white, colframe=scarlet, title={\textbf{Solution}}, breakable, grow to left by=6mm, grow to right by=10mm]
\begin{enumerate}[label=\textbf{(\alph*)}]
  \item First divide the differential equation by $\cos x$:
\[
\frac{\mathrm{d}y}{\mathrm{d}x}+y\tan x=x\sec x
\]

This is now in the linear form
\[
\frac{\mathrm{d}y}{\mathrm{d}x}+P(x)y=Q(x)
\]
with
\[
P(x)=\tan x
\]

So an integrating factor is
\begin{align*}
\mu(x)&=e^{\int P(x)\,\mathrm{d}x} \\
&=e^{\int \tan x\,\mathrm{d}x} \\
&=e^{\int \frac{\sin x}{\cos x}\,\mathrm{d}x} \\
&=e^{-\ln|\cos x|} \\
&=e^{\ln|\sec x|}
\end{align*}

Since $0<x<\frac{\pi}{2}$, we have $\cos x>0$, so $|\sec x|=\sec x$. Hence
\[
\mu(x)=e^{\ln(\sec x)}=\sec x
\]

Therefore an appropriate integrating factor is $\sec x$.

  \item Using the result from part (a), the equation is
\[
\frac{\mathrm{d}y}{\mathrm{d}x}+y\tan x=x\sec x
\]

Multiply through by the integrating factor $\sec x$:
\[
\sec x\frac{\mathrm{d}y}{\mathrm{d}x}+y\sec x\tan x=x\sec^2 x
\]

The left-hand side is the derivative of $y\sec x$, so
\[
\frac{\mathrm{d}}{\mathrm{d}x}(y\sec x)=x\sec^2 x
\]

Integrating both sides gives
\[
y\sec x=\int x\sec^2 x\,\mathrm{d}x+C
\]

Integrate by parts with
\[
u=x,\qquad \mathrm{d}v=\sec^2 x\,\mathrm{d}x
\]
so that
\[
\mathrm{d}u=\mathrm{d}x,\qquad v=\tan x
\]

Then
\begin{align*}
\int x\sec^2 x\,\mathrm{d}x
&=x\tan x-\int \tan x\,\mathrm{d}x \\
&=x\tan x-(-\ln(\cos x)) \\
&=x\tan x+\ln(\cos x)
\end{align*}

Therefore
\[
y\sec x=x\tan x+\ln(\cos x)+C
\]

Now multiply through by $\cos x$:
\begin{align*}
y&=\cos x\left(x\tan x+\ln(\cos x)+C\right) \\
&=x\sin x+\cos x\ln(\cos x)+C\cos x
\end{align*}

Use the condition $y=1$ when $x=\frac{\pi}{4}$. Substituting into
\[
y\sec x=x\tan x+\ln(\cos x)+C
\]
gives
\begin{align*}
1\cdot \sec\frac{\pi}{4}
&=\frac{\pi}{4}\tan\frac{\pi}{4}+\ln\left(\cos\frac{\pi}{4}\right)+C \\
\sqrt{2}
&=\frac{\pi}{4}+\ln\left(\frac{\sqrt{2}}{2}\right)+C \\
\sqrt{2}
&=\frac{\pi}{4}-\frac{1}{2}\ln 2+C
\end{align*}

So
\[
C=\sqrt{2}-\frac{\pi}{4}+\frac{1}{2}\ln 2
\]

Substitute this into the expression for $y$:
\[
y=x\sin x+\cos x\left[\ln(\cos x)+\sqrt{2}-\frac{\pi}{4}+\frac{1}{2}\ln 2\right]
\]

Hence the required solution is
\[
y=x\sin x+\cos x\left[\ln(\cos x)+\sqrt{2}-\frac{\pi}{4}+\frac{1}{2}\ln 2\right]
\]
  \end{enumerate}
\end{tcolorbox}


\newpage

% ---- Question 5 ----
\questionitem

\begin{enumerate}[label=\textbf{(\alph*)}]
  \item Show that, for $x>0$,
  \[
  \frac{\mathrm{d}}{\mathrm{d}x}\!\left(\ln\!\left(\coth \frac{x}{2}\right)\right)=-\operatorname{cosech} x
  \]
  \marks[-20pt]{3}
  \item Hence solve the differential equation
  \[
  \frac{\mathrm{d}y}{\mathrm{d}x}-\operatorname{cosech} x\,y=\tanh \frac{x}{2},\qquad x>0
  \]
  for which $y=1$ when $x=\ln 3$. Give your answer in exact form \marks{7} \\
\end{enumerate}

\vspace*{10pt}

\begin{tcolorbox}[colback=white, colframe=scarlet, title={\textbf{Solution}}, breakable, grow to left by=6mm, grow to right by=10mm]
\begin{enumerate}[label=\textbf{(\alph*)}]
  \item Let
\[
u=\frac{x}{2}
\]
Then, by the chain rule,
\begin{align*}
\frac{\mathrm{d}}{\mathrm{d}x}\left(\ln\left(\coth \frac{x}{2}\right)\right)
&= \frac{\mathrm{d}}{\mathrm{d}x}\bigl(\ln(\coth u)\bigr) \\
&= \frac{1}{\coth u}\cdot \frac{\mathrm{d}}{\mathrm{d}x}(\coth u) \\
&= \frac{1}{\coth u}\cdot \left(-\operatorname{cosech}^2 u\right)\cdot \frac{1}{2}
\end{align*}
Now simplify:
\begin{align*}
\frac{1}{\coth u}\cdot \left(-\operatorname{cosech}^2 u\right)\cdot \frac{1}{2}
&= -\frac{1}{2}\cdot \frac{\operatorname{cosech}^2 u}{\coth u} \\
&= -\frac{1}{2}\cdot \frac{1/\sinh^2 u}{\cosh u/\sinh u} \\
&= -\frac{1}{2\sinh u\cosh u}
\end{align*}
Using \(\sinh(2u)=2\sinh u\cosh u\),
\begin{align*}
-\frac{1}{2\sinh u\cosh u}
&= -\frac{1}{\sinh(2u)} \\
&= -\frac{1}{\sinh x} \\
&= -\operatorname{cosech} x
\end{align*}
Therefore,
\[
\frac{\mathrm{d}}{\mathrm{d}x}\left(\ln\left(\coth \frac{x}{2}\right)\right)=-\operatorname{cosech} x
\]
as required.

  \item The differential equation is
\[
\frac{\mathrm{d}y}{\mathrm{d}x}-\operatorname{cosech} x\,y=\tanh \frac{x}{2}
\]
This is linear, with integrating factor
\begin{align*}
\mu(x)
&= e^{\int -\operatorname{cosech} x\,\mathrm{d}x} \\
&= e^{\ln\left(\coth \frac{x}{2}\right)} \\
&= \coth \frac{x}{2}
\end{align*}
using part (a).

Next, from part (a),
\[
\frac{\mathrm{d}}{\mathrm{d}x}\left(\ln\left(\coth \frac{x}{2}\right)\right)=-\operatorname{cosech} x
\]
so
\begin{align*}
\frac{\mathrm{d}}{\mathrm{d}x}\left(\coth \frac{x}{2}\right)
&= \coth \frac{x}{2}\cdot \frac{\mathrm{d}}{\mathrm{d}x}\left(\ln\left(\coth \frac{x}{2}\right)\right) \\
&= -\operatorname{cosech} x\,\coth \frac{x}{2}
\end{align*}
Multiply the differential equation by \(\coth \dfrac{x}{2}\):
\begin{align*}
\coth \frac{x}{2}\,\frac{\mathrm{d}y}{\mathrm{d}x}
-\operatorname{cosech} x\,\coth \frac{x}{2}\,y
&= \coth \frac{x}{2}\tanh \frac{x}{2}
\end{align*}
The left-hand side is
\[
\frac{\mathrm{d}}{\mathrm{d}x}\left(y\coth \frac{x}{2}\right)
\]
and the right-hand side is \(1\), since \(\coth \dfrac{x}{2}\tanh \dfrac{x}{2}=1\). Hence
\[
\frac{\mathrm{d}}{\mathrm{d}x}\left(y\coth \frac{x}{2}\right)=1
\]
Integrating,
\[
y\coth \frac{x}{2}=\int 1\,\mathrm{d}x=x+C
\]
So
\[
y=(x+C)\tanh \frac{x}{2}
\]
Now use \(y=1\) when \(x=\ln 3\):
\[
1=(\ln 3+C)\tanh \frac{\ln 3}{2}
\]
Evaluate \(\tanh \dfrac{\ln 3}{2}\):
\begin{align*}
\tanh \frac{\ln 3}{2}
&= \frac{e^{\ln 3/2}-e^{-\ln 3/2}}{e^{\ln 3/2}+e^{-\ln 3/2}} \\
&= \frac{e^{\ln 3}-1}{e^{\ln 3}+1} \\
&= \frac{3-1}{3+1} \\
&= \frac{1}{2}
\end{align*}
Therefore,
\begin{align*}
1&=(\ln 3+C)\cdot \frac{1}{2} \\
2&=\ln 3+C \\
C&=2-\ln 3
\end{align*}
Substituting back gives
\[
y=(x+2-\ln 3)\tanh \frac{x}{2}
\]
Therefore, the required exact solution is \(y=(x+2-\ln 3)\tanh \dfrac{x}{2}\).
\end{enumerate}
\end{tcolorbox}


\newpage

% ---- Question 6 ----
\questionitem

\begin{enumerate}[label=\textbf{(\alph*)}]
  \item Show that an appropriate integrating factor for
  \[
  x(x-1)\,\frac{\mathrm{d}y}{\mathrm{d}x}+(x+1)y=x(x-1)^2
  \]
  is $\dfrac{(x-1)^2}{x}$ \marks{4} \\
  \item Hence find the solution of the differential equation
  \[
  x(x-1)\,\frac{\mathrm{d}y}{\mathrm{d}x}+(x+1)y=x(x-1)^2
  \]
  for $x>1$, given that $y=2$ when $x=2$ \marks{6} \\
\end{enumerate}

\vspace*{10pt}

\begin{tcolorbox}[colback=white, colframe=scarlet, title={\textbf{Solution}}, breakable, grow to left by=6mm, grow to right by=10mm]
\begin{enumerate}[label=\textbf{(\alph*)}]
  \item First write the differential equation in linear form by dividing through by $x(x-1)$:
\[
\frac{\mathrm{d}y}{\mathrm{d}x}+\frac{x+1}{x(x-1)}y=x-1
\]

For a linear equation
\[
\frac{\mathrm{d}y}{\mathrm{d}x}+P(x)y=Q(x),
\]
an integrating factor is
\[
\mu(x)=e^{\int P(x)\,\mathrm{d}x}
\]

Here
\[
P(x)=\frac{x+1}{x(x-1)}
\]

Decompose into partial fractions:
\begin{align*}
\frac{x+1}{x(x-1)}&=\frac{A}{x}+\frac{B}{x-1} \\
x+1&=A(x-1)+Bx \\
x+1&=(A+B)x-A
\end{align*}

Equating coefficients gives
\[
A+B=1,\qquad -A=1
\]
so
\[
A=-1,\qquad B=2
\]

Hence
\[
\frac{x+1}{x(x-1)}=-\frac{1}{x}+\frac{2}{x-1}
\]

Therefore
\begin{align*}
\mu(x)&=e^{\int\left(-\frac{1}{x}+\frac{2}{x-1}\right)\,\mathrm{d}x} \\
&=e^{-\ln x+2\ln(x-1)} \\
&=\frac{(x-1)^2}{x}
\end{align*}

Thus an appropriate integrating factor is
\[
\mu(x)=\frac{(x-1)^2}{x}
\]

  \item Using the integrating factor from part (a),
\[
\mu(x)=\frac{(x-1)^2}{x}
\]

Multiply the linear form
\[
\frac{\mathrm{d}y}{\mathrm{d}x}+\frac{x+1}{x(x-1)}y=x-1
\]
through by $\mu(x)$:
\[
\frac{(x-1)^2}{x}\frac{\mathrm{d}y}{\mathrm{d}x}+\frac{(x-1)^2}{x}\cdot\frac{x+1}{x(x-1)}\,y=\frac{(x-1)^2}{x}(x-1)
\]

Since $\mu$ is an integrating factor, the left-hand side becomes
\[
\frac{\mathrm{d}}{\mathrm{d}x}\left(\frac{(x-1)^2}{x}y\right)
\]

So
\[
\frac{\mathrm{d}}{\mathrm{d}x}\left(\frac{(x-1)^2}{x}y\right)=\frac{(x-1)^3}{x}
\]

Integrating both sides:
\[
\frac{(x-1)^2}{x}y=\int \frac{(x-1)^3}{x}\,\mathrm{d}x
\]

Expand the integrand:
\begin{align*}
\frac{(x-1)^3}{x}&=\frac{x^3-3x^2+3x-1}{x} \\
&=x^2-3x+3-\frac{1}{x}
\end{align*}

Hence
\begin{align*}
\frac{(x-1)^2}{x}y&=\int \left(x^2-3x+3-\frac{1}{x}\right)\,\mathrm{d}x \\
&=\frac{x^3}{3}-\frac{3x^2}{2}+3x-\ln x+C
\end{align*}

So
\[
y=\frac{x}{(x-1)^2}\left(\frac{x^3}{3}-\frac{3x^2}{2}+3x-\ln x+C\right)
\]

Now use the condition $y=2$ when $x=2$:
\begin{align*}
2&=\frac{2}{(2-1)^2}\left(\frac{2^3}{3}-\frac{3(2^2)}{2}+3(2)-\ln 2+C\right) \\
2&=2\left(\frac{8}{3}-6+6-\ln 2+C\right) \\
2&=2\left(\frac{8}{3}-\ln 2+C\right)
\end{align*}

Divide by $2$:
\[
1=\frac{8}{3}-\ln 2+C
\]

So
\[
C=1-\frac{8}{3}+\ln 2=\ln 2-\frac{5}{3}
\]

Substituting back:
\[
y=\frac{x}{(x-1)^2}\left(\frac{x^3}{3}-\frac{3x^2}{2}+3x-\ln x+\ln 2-\frac{5}{3}\right)
\]

Therefore the required solution is
\[
y=\frac{x}{(x-1)^2}\left(\frac{x^3}{3}-\frac{3x^2}{2}+3x-\ln x+\ln 2-\frac{5}{3}\right)
\]
\end{enumerate}
\end{tcolorbox}


\newpage

% ---- Question 7 ----
\questionitem

Find the general solution of the differential equation\\
\[
\frac{\mathrm{d}y}{\mathrm{d}x}+\cot x\,y=\frac{\cos x}{\sin^2 x\left(1+\ln^2(\sin x)\right)}
\]\\
where $0<x<\pi$ \marks{7}

\vspace*{10pt}

\begin{tcolorbox}[colback=white, colframe=scarlet, title={\textbf{Solution}}, breakable, grow to left by=6mm, grow to right by=10mm]
This is a first-order linear differential equation:
\[
\frac{\mathrm{d}y}{\mathrm{d}x}+\cot x\,y=\frac{\cos x}{\sin^2 x\left(1+\ln^2(\sin x)\right)}
\]

Its integrating factor is
\begin{align*}
\mu(x)
&=e^{\int \cot x\,\mathrm{d}x} \\
&=e^{\int \frac{\cos x}{\sin x}\,\mathrm{d}x} \\
&=e^{\ln(\sin x)}
\end{align*}

Since $0<x<\pi$, we have $\sin x>0$, so
\[
\mu(x)=\sin x
\]

Multiply the whole differential equation by $\sin x$:
\begin{align*}
\sin x\frac{\mathrm{d}y}{\mathrm{d}x}+\sin x\cot x\,y
&=\sin x\cdot \frac{\cos x}{\sin^2 x\left(1+\ln^2(\sin x)\right)} \\
\sin x\frac{\mathrm{d}y}{\mathrm{d}x}+\cos x\,y
&=\frac{\cos x}{\sin x\left(1+\ln^2(\sin x)\right)}
\end{align*}

Now
\[
\frac{\mathrm{d}}{\mathrm{d}x}(y\sin x)=\sin x\frac{\mathrm{d}y}{\mathrm{d}x}+\cos x\,y
\]

So the equation becomes
\[
\frac{\mathrm{d}}{\mathrm{d}x}(y\sin x)=\frac{\cos x}{\sin x\left(1+\ln^2(\sin x)\right)}
\]

Integrating both sides gives
\[
\int \frac{\mathrm{d}}{\mathrm{d}x}(y\sin x)\,\,\mathrm{d}x
=
\int \frac{\cos x}{\sin x\left(1+\ln^2(\sin x)\right)}\,\,\mathrm{d}x
\]

For the right-hand side, let
\[
u=\ln(\sin x)
\]

Then
\begin{align*}
\frac{\mathrm{d}u}{\mathrm{d}x}
&=\frac{1}{\sin x}\cdot \cos x \\
&=\frac{\cos x}{\sin x}
\end{align*}

Hence
\[
\mathrm{d}u=\frac{\cos x}{\sin x}\,\mathrm{d}x
\]

So
\begin{align*}
\int \frac{\cos x}{\sin x\left(1+\ln^2(\sin x)\right)}\,\,\mathrm{d}x
&=\int \frac{1}{1+u^2}\,\mathrm{d}u \\
&=\arctan(u)+C \\
&=\arctan(\ln(\sin x))+C
\end{align*}

Therefore
\[
y\sin x=\arctan(\ln(\sin x))+C
\]

Finally, divide by $\sin x$:
\[
y=\frac{\arctan(\ln(\sin x))+C}{\sin x}
\]

The general solution is $y=\dfrac{\arctan(\ln(\sin x))+C}{\sin x}$.
\end{tcolorbox}


\newpage

% ---- Question 8 ----
\questionitem

\begin{enumerate}[label=\textbf{(\alph*)}]
  \item Solve the differential equation
  \[
  (1+x)\,\frac{\mathrm{d}y}{\mathrm{d}x}-2y=(1+x)^4,\qquad x>-1
  \]
  giving your answer in the form $y=f(x)$ \marks{3} \\
  \item The solution curve passes through the point $(0,-1)$. Find the exact positive value of $x$ for which $y=0$ \marks{3} \\
\end{enumerate}

\vspace*{10pt}

\begin{tcolorbox}[colback=white, colframe=scarlet, title={\textbf{Solution}}, breakable, grow to left by=6mm, grow to right by=10mm]
\begin{enumerate}[label=\textbf{(\alph*)}]
  \item Rewrite the differential equation in linear form by dividing through by $1+x$:
\[
\frac{\mathrm{d}y}{\mathrm{d}x}-\frac{2}{1+x}y=(1+x)^3
\]

Since this is a linear equation, the integrating factor is
\begin{align*}
\mu(x)&=e^{\int -\frac{2}{1+x}\,\,\mathrm{d}x} \\
&=e^{-2\ln(1+x)} \\
&=(1+x)^{-2}
\end{align*}

Multiplying the differential equation by $(1+x)^{-2}$ gives
\begin{align*}
(1+x)^{-2}\frac{\mathrm{d}y}{\mathrm{d}x}-2(1+x)^{-3}y&=(1+x)
\end{align*}

The left-hand side is the derivative of $y(1+x)^{-2}$, so
\[
\frac{\mathrm{d}}{\mathrm{d}x}\bigl(y(1+x)^{-2}\bigr)=1+x
\]

Integrating both sides:
\begin{align*}
y(1+x)^{-2}&=\int (1+x)\,\,\mathrm{d}x \\
&=x+\frac{x^2}{2}+C
\end{align*}

Hence
\[
y=(1+x)^2\left(\frac{x^2}{2}+x+C\right)
\]

Now use the point $(0,-1)$:
\begin{align*}
-1&=(1+0)^2\left(\frac{0^2}{2}+0+C\right) \\
-1&=C
\end{align*}

Therefore
\[
y=(1+x)^2\left(\frac{x^2}{2}+x-1\right)
\]

So the required function is $y=(1+x)^2\left(\frac{x^2}{2}+x-1\right)$.

  \item To find when $y=0$, set the expression from part (a) equal to zero:
\[
(1+x)^2\left(\frac{x^2}{2}+x-1\right)=0
\]

Since $x>-1$, we have $1+x>0$, so $(1+x)^2\neq 0$. Therefore
\[
\frac{x^2}{2}+x-1=0
\]

Multiply by $2$:
\[
x^2+2x-2=0
\]

Solve using the quadratic formula:
\begin{align*}
x&=\frac{-2\pm\sqrt{2^2-4(1)(-2)}}{2} \\
&=\frac{-2\pm\sqrt{4+8}}{2} \\
&=\frac{-2\pm\sqrt{12}}{2} \\
&=\frac{-2\pm 2\sqrt{3}}{2} \\
&=-1\pm \sqrt{3}
\end{align*}

The positive value is
\[
x=\sqrt{3}-1
\]

So the exact positive value of $x$ is $x=\sqrt{3}-1$.
\end{enumerate}
\end{tcolorbox}


\newpage

% ---- Question 9 ----
\questionitem

Find the general solution of the differential equation\\
\[
\frac{\mathrm{d}y}{\mathrm{d}x}-2y\tan x=\sec^3 x
\]\\
where $-\frac{\pi}{2}<x<\frac{\pi}{2}$ \marks{7}

\vspace*{10pt}

\begin{tcolorbox}[colback=white, colframe=scarlet, title={\textbf{Solution}}, breakable, grow to left by=6mm, grow to right by=10mm]
Write the differential equation in the linear form
\[
\frac{\mathrm{d}y}{\mathrm{d}x}+P(x)y=Q(x)
\]
so here
\[
\frac{\mathrm{d}y}{\mathrm{d}x}-2y\tan x=\sec^3 x,\qquad P(x)=-2\tan x,\qquad Q(x)=\sec^3 x
\]

For a linear equation, the integrating factor is
\[
\mu(x)=e^{\int P(x)\,\,\mathrm{d}x}
\]

So
\begin{align*}
\mu(x)&=e^{\int -2\tan x\,\,\mathrm{d}x} \\
&=e^{-2\int \frac{\sin x}{\cos x}\,\,\mathrm{d}x} \\
&=e^{-2(-\ln(\cos x))} \\
&=e^{2\ln(\cos x)} \\
&=\cos^2 x
\end{align*}

Since $-\frac{\pi}{2}<x<\frac{\pi}{2}$, we have $\cos x>0$, so this is valid without absolute values.

Now multiply the whole differential equation by $\cos^2 x$:
\begin{align*}
\cos^2 x\frac{\mathrm{d}y}{\mathrm{d}x}-2y\cos^2 x\tan x&=\cos^2 x\sec^3 x \\
\cos^2 x\frac{\mathrm{d}y}{\mathrm{d}x}-2y\cos x\sin x&=\sec x
\end{align*}

Next, observe that
\begin{align*}
\frac{\mathrm{d}}{\mathrm{d}x}(y\cos^2 x)
&=\cos^2 x\frac{\mathrm{d}y}{\mathrm{d}x}+y\frac{\mathrm{d}}{\mathrm{d}x}(\cos^2 x) \\
&=\cos^2 x\frac{\mathrm{d}y}{\mathrm{d}x}+y(2\cos x)(-\sin x) \\
&=\cos^2 x\frac{\mathrm{d}y}{\mathrm{d}x}-2y\cos x\sin x
\end{align*}

So the left-hand side is an exact derivative, and the equation becomes
\[
\frac{\mathrm{d}}{\mathrm{d}x}(y\cos^2 x)=\sec x
\]

Integrating both sides with respect to $x$ gives
\[
y\cos^2 x=\int \sec x\,\,\mathrm{d}x
\]

To integrate $\sec x$, use the standard trick:
\begin{align*}
\int \sec x\,\,\mathrm{d}x
&=\int \sec x\cdot \frac{\sec x+\tan x}{\sec x+\tan x}\,\,\mathrm{d}x \\
&=\int \frac{\sec^2 x+\sec x\tan x}{\sec x+\tan x}\,\,\mathrm{d}x
\end{align*}

Let
\[
u=\sec x+\tan x
\]
Then
\[
\frac{\mathrm{d}u}{\mathrm{d}x}=\sec x\tan x+\sec^2 x
\]
so
\[
\int \sec x\,\,\mathrm{d}x=\int \frac{1}{u}\,\,\mathrm{d}u=\ln|u|+C=\ln|\sec x+\tan x|+C
\]

On the interval $-\frac{\pi}{2}<x<\frac{\pi}{2}$, we have
\[
\sec x+\tan x=\frac{1+\sin x}{\cos x}>0
\]
because $1+\sin x>0$ and $\cos x>0$. Hence
\[
\ln|\sec x+\tan x|=\ln(\sec x+\tan x)
\]

Therefore
\[
y\cos^2 x=\ln(\sec x+\tan x)+C
\]

Finally, divide through by $\cos^2 x$:
\[
y=\sec^2 x\bigl(\ln(\sec x+\tan x)+C\bigr)
\]

Hence the general solution is
\[
y=\sec^2 x\bigl(\ln(\sec x+\tan x)+C\bigr)
\]
\end{tcolorbox}


\newpage

% ---- Question 10 ----
\questionitem

\begin{enumerate}[label=\textbf{(\alph*)}]
  \item Show that
  \[
  \frac{\mathrm{d}}{\mathrm{d}x}\!\left(\ln\!\left(\frac{\cosh x+1}{\sinh x}\right)\right)=-\operatorname{cosech} x
  \]
  \marks[-20pt]{3}
  \item Hence solve, for $x>0$,
  \[
  \sinh x\,\frac{\mathrm{d}y}{\mathrm{d}x}-y=1-\cosh x
  \]
  given that $y=2$ when $x=\ln 3$. Give your answer in exact form \marks{7} \\
\end{enumerate}

\vspace*{10pt}

\begin{tcolorbox}[colback=white, colframe=scarlet, title={\textbf{Solution}}, breakable, grow to left by=6mm, grow to right by=10mm]
\begin{enumerate}[label=\textbf{(\alph*)}]
  \item Write
\[
\ln\left(\frac{\cosh x+1}{\sinh x}\right)=\ln(\cosh x+1)-\ln(\sinh x)
\]

Differentiate term by term:
\begin{align*}
\frac{\mathrm{d}}{\mathrm{d}x}\left(\ln\left(\frac{\cosh x+1}{\sinh x}\right)\right)
&=\frac{\sinh x}{\cosh x+1}-\frac{\cosh x}{\sinh x} \\
&=\frac{\sinh^2 x-\cosh x(\cosh x+1)}{\sinh x(\cosh x+1)}
\end{align*}

Now simplify the numerator using $\cosh^2 x-\sinh^2 x=1$:
\begin{align*}
\sinh^2 x-\cosh x(\cosh x+1)
&=\sinh^2 x-\cosh^2 x-\cosh x \\
&=-(\cosh^2 x-\sinh^2 x)-\cosh x \\
&=-1-\cosh x \\
&=-(\cosh x+1)
\end{align*}

So
\begin{align*}
\frac{\mathrm{d}}{\mathrm{d}x}\left(\ln\left(\frac{\cosh x+1}{\sinh x}\right)\right)
&=\frac{-(\cosh x+1)}{\sinh x(\cosh x+1)} \\
&=-\frac{1}{\sinh x} \\
&=-\operatorname{cosech}x
\end{align*}

Hence,
\[
\frac{\mathrm{d}}{\mathrm{d}x}\left(\ln\left(\frac{\cosh x+1}{\sinh x}\right)\right)=-\operatorname{cosech}x
\]

  \item Start with
\[
\sinh x\,\frac{\mathrm{d}y}{\mathrm{d}x}-y=1-\cosh x
\]

Since $x>0$, we have $\sinh x\neq 0$, so divide through by $\sinh x$:
\[
\frac{\mathrm{d}y}{\mathrm{d}x}-\operatorname{cosech}x\,y=\frac{1-\cosh x}{\sinh x}
\]

This is a linear differential equation of the form
\[
\frac{\mathrm{d}y}{\mathrm{d}x}+P(x)y=Q(x)
\]
with
\[
P(x)=-\operatorname{cosech}x
\]

From part (a),
\[
\frac{\mathrm{d}}{\mathrm{d}x}\left(\ln\left(\frac{\cosh x+1}{\sinh x}\right)\right)=-\operatorname{cosech}x
\]
so an integrating factor is
\[
\mu(x)=\exp\left(\int -\operatorname{cosech}x\,\mathrm{d}x\right)=\frac{\cosh x+1}{\sinh x}
\]

Multiply the differential equation throughout by $\mu(x)$:
\[
\frac{\cosh x+1}{\sinh x}\frac{\mathrm{d}y}{\mathrm{d}x}-\frac{\cosh x+1}{\sinh^2 x}y
=\frac{(\cosh x+1)(1-\cosh x)}{\sinh^2 x}
\]

The left-hand side is
\[
\frac{\mathrm{d}}{\mathrm{d}x}\left(\mu y\right)
=\frac{\mathrm{d}}{\mathrm{d}x}\left(\frac{\cosh x+1}{\sinh x}y\right)
\]

Simplify the right-hand side:
\begin{align*}
\frac{(\cosh x+1)(1-\cosh x)}{\sinh^2 x}
&=\frac{1-\cosh^2 x}{\sinh^2 x} \\
&=\frac{-\sinh^2 x}{\sinh^2 x} \\
&=-1
\end{align*}

So
\[
\frac{\mathrm{d}}{\mathrm{d}x}\left(\frac{\cosh x+1}{\sinh x}y\right)=-1
\]

Integrate with respect to $x$:
\[
\frac{\cosh x+1}{\sinh x}y=-x+C
\]

Hence
\[
y=\frac{(C-x)\sinh x}{\cosh x+1}
\]

Now use the condition $y=2$ when $x=\ln 3$.

First find $\sinh(\ln 3)$ and $\cosh(\ln 3)$:
\begin{align*}
\sinh(\ln 3)&=\frac{e^{\ln 3}-e^{-\ln 3}}{2}=\frac{3-\frac13}{2}=\frac{4}{3} \\
\cosh(\ln 3)&=\frac{e^{\ln 3}+e^{-\ln 3}}{2}=\frac{3+\frac13}{2}=\frac{5}{3}
\end{align*}

Substitute into the solution:
\begin{align*}
2
&=\frac{(C-\ln 3)\sinh(\ln 3)}{\cosh(\ln 3)+1} \\
&=\frac{(C-\ln 3)\cdot \frac43}{\frac53+1} \\
&=\frac{(C-\ln 3)\cdot \frac43}{\frac83} \\
&=\frac{C-\ln 3}{2}
\end{align*}

Therefore
\begin{align*}
2&=\frac{C-\ln 3}{2} \\
4&=C-\ln 3 \\
C&=4+\ln 3
\end{align*}

Substituting back,
\[
y=\frac{(4+\ln 3-x)\sinh x}{\cosh x+1}
\]

So the exact solution is
\[
y=\frac{(4+\ln 3-x)\sinh x}{\cosh x+1}
\]
\end{enumerate}
\end{tcolorbox}


\newpage

% ---- Question 11 ----
\questionitem

\begin{enumerate}[label=\textbf{(\alph*)}]
  \item Find the general solution of the differential equation\\
  \[
  \frac{\mathrm{d}y}{\mathrm{d}x}+\frac{2x}{x^2+1}y=6x
  \]
  giving your answer in the form $y=f(x)$ \marks{3} \\
  \item Given that $y=\tfrac{7}{2}$ when $x=0$, find the smallest positive value of $x$ for which $y=4$ \marks{3} \\
\end{enumerate}

\vspace*{10pt}

\begin{tcolorbox}[colback=white, colframe=scarlet, title={\textbf{Solution}}, breakable, grow to left by=6mm, grow to right by=10mm]
\begin{enumerate}[label=\textbf{(\alph*)}]
  \item The differential equation is linear:
  \[
  \frac{\mathrm{d}y}{\mathrm{d}x}+\frac{2x}{x^2+1}y=6x
  \]

  Its integrating factor is
  \begin{align*}
  \mu(x)&=e^{\int \frac{2x}{x^2+1}\,\,\mathrm{d}x} \\
  &=e^{\ln(x^2+1)} \\
  &=x^2+1
  \end{align*}

  Multiplying the differential equation by $x^2+1$ gives
  \begin{align*}
  (x^2+1)\frac{\mathrm{d}y}{\mathrm{d}x}+2xy&=6x(x^2+1)
  \end{align*}

  The left-hand side is the derivative of $(x^2+1)y$, so
  \[
  \frac{\mathrm{d}}{\mathrm{d}x}\bigl((x^2+1)y\bigr)=6x(x^2+1)
  \]

  Integrating both sides:
  \begin{align*}
  (x^2+1)y&=\int 6x(x^2+1)\,\,\mathrm{d}x \\
  &=\int (6x^3+6x)\,\,\mathrm{d}x \\
  &=\frac{3}{2}x^4+3x^2+C
  \end{align*}

  Hence
  \[
  y=\frac{\frac{3}{2}x^4+3x^2+C}{x^2+1}
  \]

  Therefore, the general solution is
  \[
  y=\frac{\frac{3}{2}x^4+3x^2+C}{x^2+1}
  \]

  \item Using $y=\tfrac{7}{2}$ when $x=0$:
  \begin{align*}
  \frac{7}{2}&=\frac{\frac{3}{2}(0)^4+3(0)^2+C}{(0)^2+1} \\
  \frac{7}{2}&=C
  \end{align*}

  So the particular solution is
  \[
  y=\frac{\frac{3}{2}x^4+3x^2+\frac{7}{2}}{x^2+1}
  \]

  Now set $y=4$:
  \begin{align*}
  4&=\frac{\frac{3}{2}x^4+3x^2+\frac{7}{2}}{x^2+1}
  \end{align*}

  Multiply through by $x^2+1$:
  \begin{align*}
  4(x^2+1)&=\frac{3}{2}x^4+3x^2+\frac{7}{2}
  \end{align*}

  Multiply by $2$:
  \begin{align*}
  8x^2+8&=3x^4+6x^2+7
  \end{align*}

  Rearranging:
  \begin{align*}
  0&=3x^4-2x^2-1
  \end{align*}

  Let $u=x^2$. Then
  \begin{align*}
  3u^2-2u-1&=0 \\
  (3u+1)(u-1)&=0
  \end{align*}

  So
  \[
  u=1 \quad \text{or} \quad u=-\frac{1}{3}
  \]

  Since $u=x^2\geq 0$, we take $x^2=1$, giving
  \[
  x=\pm 1
  \]

  The smallest positive value is
  \[
  x=1
  \]
\end{enumerate}
\end{tcolorbox}


\newpage

% ---- Question 12 ----
\questionitem

For $x>0$, solve the differential equation\\
\[
\sinh x\,\frac{\mathrm{d}y}{\mathrm{d}x} - 2y\cosh x = \sinh^2 x\cosh x
\] \\
given that $y=5$ when $x=\ln(1+\sqrt{2})$. \\
Give your answer in exact form. \marks{7}

\vspace*{10pt}

\begin{tcolorbox}[colback=white, colframe=scarlet, title={\textbf{Solution}}, breakable, grow to left by=6mm, grow to right by=10mm]
Since $x>0$, we have $\sinh x>0$, so in particular $\sinh x \neq 0$. We may therefore divide the differential equation by $\sinh x$ to put it into linear form:
\begin{align*}
\sinh x\,\frac{\mathrm{d}y}{\mathrm{d}x} - 2y\cosh x &= \sinh^2 x\cosh x \\
\frac{\mathrm{d}y}{\mathrm{d}x} - 2\frac{\cosh x}{\sinh x}y &= \sinh x\cosh x \\
\frac{\mathrm{d}y}{\mathrm{d}x} - 2\coth x\,y &= \sinh x\cosh x
\end{align*}

This is a first-order linear equation of the form
\[
\frac{\mathrm{d}y}{\mathrm{d}x} + P(x)y = Q(x)
\]
with
\[
P(x)=-2\coth x
\]

So the integrating factor is
\begin{align*}
\mu(x) &= e^{\int -2\coth x\,\mathrm{d}x} \\
&= e^{-2\int \coth x\,\mathrm{d}x} \\
&= e^{-2\ln(\sinh x)} \\
&= \frac{1}{\sinh^2 x}
\end{align*}

Multiplying the differential equation by this integrating factor gives
\begin{align*}
\frac{1}{\sinh^2 x}\frac{\mathrm{d}y}{\mathrm{d}x} - \frac{2\coth x}{\sinh^2 x}y
&= \frac{\sinh x\cosh x}{\sinh^2 x} \\
\frac{\mathrm{d}}{\mathrm{d}x}\left(\frac{y}{\sinh^2 x}\right) &= \coth x
\end{align*}

Now integrate both sides with respect to $x$:
\begin{align*}
\frac{y}{\sinh^2 x} &= \int \coth x\,\mathrm{d}x \\
&= \ln(\sinh x) + C
\end{align*}

Hence
\[
y=\sinh^2 x\bigl(\ln(\sinh x)+C\bigr)
\]

Now use the condition $y=5$ when $x=\ln(1+\sqrt{2})$.

We need $\sinh\!\bigl(\ln(1+\sqrt{2})\bigr)$. Using
\[
\sinh t=\frac{e^t-e^{-t}}{2}
\]
with $t=\ln(1+\sqrt{2})$, we get
\begin{align*}
\sinh\!\bigl(\ln(1+\sqrt{2})\bigr)
&= \frac{(1+\sqrt{2})-\dfrac{1}{1+\sqrt{2}}}{2} \\
&= \frac{(1+\sqrt{2})-(\sqrt{2}-1)}{2} \\
&= \frac{2}{2} \\
&= 1
\end{align*}

So substituting into the solution:
\begin{align*}
5 &= 1^2\bigl(\ln 1 + C\bigr) \\
5 &= 0 + C \\
C &= 5
\end{align*}

Therefore the required exact solution is
\[
y=\sinh^2 x\bigl(\ln(\sinh x)+5\bigr)
\]

So the answer is $y=\sinh^2 x\bigl(\ln(\sinh x)+5\bigr)$.
\end{tcolorbox}


\newpage

% ---- Question 13 ----
\questionitem

For $x>0$, a function $y$ satisfies \\
\[
\sinh x\,\frac{\mathrm{d}y}{\mathrm{d}x}-y\cosh x=\sinh x\,\cosh^3 x
\] \\
\begin{enumerate}[label=\textbf{(\alph*)}]
  \item By first writing the equation in the form $\frac{\mathrm{d}y}{\mathrm{d}x}+P(x)y=Q(x)$, find an integrating factor \marks{2} \\
  \item Hence solve the differential equation, given that $y=2$ when $x=\ln(1+\sqrt{2})$ \marks{5} \\
\end{enumerate}
Give your answer in exact form.

\vspace*{10pt}

\begin{tcolorbox}[colback=white, colframe=scarlet, title={\textbf{Solution}}, breakable, grow to left by=6mm, grow to right by=10mm]
\begin{enumerate}[label=\textbf{(\alph*)}]
  \item First divide the differential equation by \(\sinh x\):
\[
\frac{\mathrm{d}y}{\mathrm{d}x}-\frac{\cosh x}{\sinh x}y=\cosh^3 x
\]
So it is in the form \(\dfrac{\mathrm{d}y}{\mathrm{d}x}+P(x)y=Q(x)\) with
\[
P(x)=-\coth x
\]

The integrating factor is
\begin{align*}
\mu(x)
&=e^{\int P(x)\,\,\mathrm{d}x} \\
&=e^{\int -\coth x\,\,\mathrm{d}x} \\
&=e^{-\ln(\sinh x)} \\
&=\frac{1}{\sinh x}
\end{align*}
since \(x>0\), so \(\sinh x>0\).

Final answer: the integrating factor is \(\dfrac{1}{\sinh x}\).

  \item Multiply the equation
\[
\frac{\mathrm{d}y}{\mathrm{d}x}-\coth x\,y=\cosh^3 x
\]
by the integrating factor \(\dfrac{1}{\sinh x}\):
\[
\frac{1}{\sinh x}\frac{\mathrm{d}y}{\mathrm{d}x}-\frac{\cosh x}{\sinh^2 x}y=\frac{\cosh^3 x}{\sinh x}
\]

Now
\[
\frac{\mathrm{d}}{\mathrm{d}x}\left(\frac{y}{\sinh x}\right)
=\frac{1}{\sinh x}\frac{\mathrm{d}y}{\mathrm{d}x}-\frac{\cosh x}{\sinh^2 x}y
\]
so the equation becomes
\[
\frac{\mathrm{d}}{\mathrm{d}x}\left(\frac{y}{\sinh x}\right)=\frac{\cosh^3 x}{\sinh x}
\]

Integrating both sides gives
\[
\frac{y}{\sinh x}=\int \frac{\cosh^3 x}{\sinh x}\,\,\mathrm{d}x+C
\]

Rewrite the integrand using \(\cosh^2 x=1+\sinh^2 x\):
\begin{align*}
\frac{\cosh^3 x}{\sinh x}
&=\frac{\cosh x\,\cosh^2 x}{\sinh x} \\
&=\frac{\cosh x(1+\sinh^2 x)}{\sinh x}
\end{align*}

Let
\[
u=\sinh x
\quad\Rightarrow\quad
\mathrm{d}u=\cosh x\,\mathrm{d}x
\]
Then
\begin{align*}
\int \frac{\cosh^3 x}{\sinh x}\,\,\mathrm{d}x
&=\int \frac{\cosh x(1+\sinh^2 x)}{\sinh x}\,\,\mathrm{d}x \\
&=\int \left(\frac{1}{u}+u\right)\,\mathrm{d}u \\
&=\ln u+\frac{u^2}{2}+C \\
&=\ln(\sinh x)+\frac{1}{2}\sinh^2 x+C
\end{align*}

Hence
\[
\frac{y}{\sinh x}=\ln(\sinh x)+\frac{1}{2}\sinh^2 x+C
\]
and so
\[
y=\sinh x\ln(\sinh x)+\frac{1}{2}\sinh^3 x+C\sinh x
\]

Use the condition \(y=2\) when \(x=\ln(1+\sqrt{2})\). First find \(\sinh\bigl(\ln(1+\sqrt{2})\bigr)\):
\begin{align*}
\sinh\bigl(\ln(1+\sqrt{2})\bigr)
&=\frac{e^{\ln(1+\sqrt{2})}-e^{-\ln(1+\sqrt{2})}}{2} \\
&=\frac{(1+\sqrt{2})-\dfrac{1}{1+\sqrt{2}}}{2} \\
&=\frac{(1+\sqrt{2})-(\sqrt{2}-1)}{2} \\
&=1
\end{align*}

So substituting \(\sinh x=1\) and \(y=2\):
\begin{align*}
2&=1\cdot \ln 1+\frac{1}{2}(1)^3+C(1) \\
2&=0+\frac{1}{2}+C \\
C&=\frac{3}{2}
\end{align*}

Therefore
\[
y=\sinh x\ln(\sinh x)+\frac{1}{2}\sinh^3 x+\frac{3}{2}\sinh x
\]

Final answer: \(y=\sinh x\ln(\sinh x)+\dfrac{1}{2}\sinh^3 x+\dfrac{3}{2}\sinh x\).
\end{enumerate}
\end{tcolorbox}


\newpage

% ---- Question 14 ----
\questionitem

For $x>0$, solve the differential equation\\
\[
\sinh x\,\frac{\mathrm{d}y}{\mathrm{d}x} - y\cosh x = \sinh x\cosh x
\] \\
given that $y = 2$ when $x = \ln(1+\sqrt{2})$. \\
Give your answer in exact form. \marks{7}

\vspace*{10pt}

\begin{tcolorbox}[colback=white, colframe=scarlet, title={\textbf{Solution}}, breakable, grow to left by=6mm, grow to right by=10mm]
Since $x>0$, we have $\sinh x>0$, so we may divide through by $\sinh x$ to put the differential equation into linear form:

\begin{align*}
\sinh x\,\frac{\mathrm{d}y}{\mathrm{d}x} - y\cosh x &= \sinh x\cosh x \\
\frac{\mathrm{d}y}{\mathrm{d}x} - \frac{\cosh x}{\sinh x}\,y &= \cosh x \\
\frac{\mathrm{d}y}{\mathrm{d}x} - \coth x\,y &= \cosh x
\end{align*}

This is a first-order linear equation of the form

\[
\frac{\mathrm{d}y}{\mathrm{d}x} + P(x)y = Q(x)
\]

with

\[
P(x) = -\coth x
\]

So the integrating factor is

\begin{align*}
\mu(x) &= e^{\int -\coth x\,\mathrm{d}x} \\
&= e^{-\int \coth x\,\mathrm{d}x} \\
&= e^{-\ln(\sinh x)} \\
&= \frac{1}{\sinh x}
\end{align*}

Multiplying the differential equation by this integrating factor gives

\begin{align*}
\frac{1}{\sinh x}\frac{\mathrm{d}y}{\mathrm{d}x} - \frac{\coth x}{\sinh x}y &= \frac{\cosh x}{\sinh x}
\end{align*}

The left-hand side is the derivative of $\dfrac{y}{\sinh x}$, so

\[
\frac{\mathrm{d}}{\mathrm{d}x}\left(\frac{y}{\sinh x}\right)=\coth x
\]

Integrating both sides with respect to $x$:

\begin{align*}
\frac{y}{\sinh x} &= \int \coth x\,\mathrm{d}x + C \\
&= \ln(\sinh x) + C
\end{align*}

Hence

\[
y=\sinh x\bigl(\ln(\sinh x)+C\bigr)
\]

Now use the condition $y=2$ when $x=\ln(1+\sqrt{2})$.

We need $\sinh\bigl(\ln(1+\sqrt{2})\bigr)$. Using
\[
\sinh u=\frac{e^u-e^{-u}}{2}
\]
with $u=\ln(1+\sqrt{2})$:

\begin{align*}
\sinh\bigl(\ln(1+\sqrt{2})\bigr)
&= \frac{(1+\sqrt{2})-\dfrac{1}{1+\sqrt{2}}}{2}
\end{align*}

Now rationalise:

\begin{align*}
\frac{1}{1+\sqrt{2}}
&= \frac{1-\sqrt{2}}{(1+\sqrt{2})(1-\sqrt{2})} \\
&= \frac{1-\sqrt{2}}{1-2} \\
&= \sqrt{2}-1
\end{align*}

So

\begin{align*}
\sinh\bigl(\ln(1+\sqrt{2})\bigr)
&= \frac{(1+\sqrt{2})-(\sqrt{2}-1)}{2} \\
&= \frac{2}{2} \\
&= 1
\end{align*}

Substituting into the general solution:

\begin{align*}
2 &= 1\bigl(\ln 1 + C\bigr) \\
2 &= 0 + C \\
C &= 2
\end{align*}

Therefore the exact solution is

\[
y=\sinh x\bigl(\ln(\sinh x)+2\bigr)
\]

So the required solution is $y=\sinh x\bigl(\ln(\sinh x)+2\bigr)$
\end{tcolorbox}


\newpage

% ---- Question 15 ----
\questionitem

Find the general solution of the differential equation\\
\[
\frac{\mathrm{d}y}{\mathrm{d}x} - y\tan x = \frac{1}{\sqrt{1+\sin x}}
\] \\
where $-\tfrac{\pi}{2} < x < \tfrac{\pi}{2}$ \marks{7}

\vspace*{10pt}

\begin{tcolorbox}[colback=white, colframe=scarlet, title={\textbf{Solution}}, breakable, grow to left by=6mm, grow to right by=10mm]
Write the differential equation in linear form:
\[
\frac{\mathrm{d}y}{\mathrm{d}x} + P(x)y = Q(x)
\]
with
\[
P(x)=-\tan x,
\qquad
Q(x)=\frac{1}{\sqrt{1+\sin x}}
\]

So the integrating factor is
\begin{align*}
\mu(x)
&= e^{\int -\tan x\,\mathrm{d}x} \\
&= e^{\ln(\cos x)}
\end{align*}

Since $-\tfrac{\pi}{2}<x<\tfrac{\pi}{2}$, we have $\cos x>0$, so
\[
\mu(x)=\cos x
\]

Multiply the differential equation throughout by $\cos x$:
\begin{align*}
\cos x\frac{\mathrm{d}y}{\mathrm{d}x} - y\cos x\tan x
&= \frac{\cos x}{\sqrt{1+\sin x}} \\
\cos x\frac{\mathrm{d}y}{\mathrm{d}x} - y\sin x
&= \frac{\cos x}{\sqrt{1+\sin x}}
\end{align*}

Now
\[
\frac{\mathrm{d}}{\mathrm{d}x}(y\cos x)=\cos x\frac{\mathrm{d}y}{\mathrm{d}x}-y\sin x
\]

So the equation becomes
\[
\frac{\mathrm{d}}{\mathrm{d}x}(y\cos x)=\frac{\cos x}{\sqrt{1+\sin x}}
\]

Integrate both sides with respect to $x$:
\[
y\cos x=\int \frac{\cos x}{\sqrt{1+\sin x}}\,\mathrm{d}x + C
\]

Let
\[
u=1+\sin x
\qquad \Rightarrow \qquad
\frac{\mathrm{d}u}{\mathrm{d}x}=\cos x
\qquad \Rightarrow \qquad
\mathrm{d}u=\cos x\,\mathrm{d}x
\]

Then
\begin{align*}
\int \frac{\cos x}{\sqrt{1+\sin x}}\,\mathrm{d}x
&= \int u^{-1/2}\,\mathrm{d}u \\
&= 2u^{1/2} \\
&= 2\sqrt{1+\sin x}
\end{align*}

Hence
\[
y\cos x=2\sqrt{1+\sin x}+C
\]

Therefore
\[
y=\frac{2\sqrt{1+\sin x}+C}{\cos x}
\]

So the general solution is
\[
y=\frac{2\sqrt{1+\sin x}+C}{\cos x}
\]

An equivalent form is
\[
y=\frac{2}{\sqrt{1-\sin x}}+C\sec x
\]
\end{tcolorbox}

\end{enumerate}
\end{document}
